\chapter{黑格尔历史哲学中的进步与辩证法}

\section*{1}

在《历史哲学讲演录》导言中，黑格尔精练地、简明扼要地将世界历史界定为“自由意识的进步”，并补充道：“【……】一种我们必须在其必然性中予以认识的进步。”因此，讲座的主题就是认识世界历史中的必然进步——进步的理论构成历史哲学的核心，自由理念则是这个理论的基准点和线索。\footnote{此文系在1969年巴黎举办的第七届黑格尔国际大会上的谈话。}除了《法哲学原理》序言的那个著名的模糊表达（“凡是现实的，都是合乎理性的；凡是合乎理性的，都是现实的”）以外，黑格尔哲学中没有第二个命题像历史哲学的这个基本命题那样，被人频繁引用，同时也被彻底误解。这两个命题的影响（同时下面也会讲到它们的联系）清楚地表明，无论是显而易见的晦涩，还是清楚明白的表述，都会使人对哲学命题产生误解和曲解。

事实上，与先前所有世纪相比，19世纪在科学技术以及与之相关的对自然的支配上都取得了无可比拟的进步，对这样一个时代来说，还有什么能比黑格尔关于自由意识的进步的历史公式更加令人信服呢？如果说，在黑格尔看来，历史哲学的任务在于认识这种进步的必然性\footnote{指前面所言及的“历史哲学的任务在于发现历史进步的必然性或规律”。——译者}，则它的追求便与19世纪实证主义学派的诉求并无二致：发现这种所有人有目共睹的历史进步背后世界历史的各个阶段或者各个时代有规则的更替规律，这种规律决定了历史事件的进程及其发展趋势。就像18世纪的杜尔哥(Turgot)和19世纪的孔德一样，黑格尔也接过该定义，提出了一种世界历史的三阶段规律，作为他的世界历史规划的基础：按照自由的规则，最初一个人是自由的，然后一些人是自由的，最后所有人都是自由的，由此，世界历史就是一个由东方世界、希腊—罗马世界和日耳曼世界或现代世界构成的时间序列。在这个自由进展的理念中，黑格尔似乎完全接近了他的时代的那种进步的乐观主义。这种乐观主义认为，所有人的自由都在其中得到实现的市民社会就是一切可能世界中最好的世界。由此，这种神正论便经历了一种莱布尼茨从未梦想过的扩张。通过认识肯定物、认识社会整体的向前发展，世界的负面物（在莱布尼茨这里是个体的恶）得到了辩护：历史的进步不是假象、不是欺骗，世界精神绝不撒谎。

然而，对黑格尔公式的误解并不在于，它似乎用一种自由的量的进展去表现世界历史的进程，而在于它以现时代的进步为尺度去打量这一进程。换言之，从最新的进展出发，将以往的历程构想为一个必然通向这一进展的发展过程。这种时间序列上的定位“预设了，在世界历史中，唯有后果丰富者方才有效，世界事件的前后序列须按成果的理性去评定”。\footnote{参见洛维特：《从黑格尔到尼采》，第3版，苏黎世，1953年，第238页。}19世纪下半叶，达尔文的进化论在自然领域中为这种成果评价提供了坚定的支持，由此，这种成果评价就构成了黑格尔的进步观念为何如此流行的原因。黑格尔的进步观使按照“强者的权利”在与封建统治势力斗争中获胜的市民阶级心安理得，就像随着工业生产力的解放，无产阶级出现为市民等级的世界历史的对手一样。然而，自由意识的进步理念却变成了市民意识的进步的意识形态。在达尔文主义进化论和实证主义社会学的影响下，历史变成了一个掩盖了历史主体自身的作用及其活动意义的线性的自然过程。\footnote{卢卡奇：《历史与阶级意识》，柏林，1923年，第215页。也参见梅洛-庞蒂：《辩证法的历险》，法兰克福，1968年，第44页以下。}进步在一定程度上抽身隐退，变成了需要依赖的事实，或是一种人们必须服从的规律。

这个偶像化了的进步概念为社会主义理论所克服，乃是19世纪历史的悖论之一。恩格斯在其著作《路德维希·费尔巴哈与德国古典哲学的终结》（1888年）中对黑格尔历史哲学的解释实现了这种突破。恩格斯与黑格尔一致，也承认自然和历史的发展有着根本的区别。自然中，只有无意识的、盲目的力量彼此发生作用，不会出现有意愿的、有意识的目的，而在社会历史中，则是有意识的、经过思虑和凭激情行动的、追求一定目的的人。但在恩格斯看来，这种区别并不能改变下述事实，即“历史过程”是受内在的普遍“规律”支配的，问题只是在于发现此种规律。\footnote{《马克思恩格斯全集》第21卷，柏林，1962年，第296页以下。（中译文参见《马克思恩格斯选集》第四卷，人民出版社2012年版，第253-254页。——译者）}这里，取代自由意识中进步理念的是一种自然历史发展的理论，这种理论根本而言是实证主义的。在这种理论中，进步变成了规律，就像支配物体在空间中的运动规律一样，它也支配了人在时间中的行为。诚如马克思所言，时间就是历史的空间。

进步思想的这种改变一直影响至今，其结果是，黑格尔的进步观念与辩证法之间的原始联系消失不见了。我的阐述旨在通过对黑格尔历史哲学中的历史运动模式及其时间结构进行解释，恢复这种联系。我可以引述一名现代经典诗人的证言，这位诗人——在同代诗人中，这是一个独一无二的现象——同时又是一个杰出的辩证法家和黑格尔哲学的爱好者。我指的是布莱希特(Bertold Brecht)。在他不久前出版的遗存下来的哲学笔记中，对社会主义者们的“莽撞的进步概念”做出了这样的简洁而又难以捉摸的表述：“‘进步’这个概念具有非常愉快的政治性质，但它对‘辩证法’的概念却有不利的后果。”\footnote{布莱希特：《政治与社会文集》(Schriften zur Politik und Gesellschaft)第1卷，W. 黑希特(Hecht)出版，柏林，1968年，第250页。}针对线性进步导致了相反方向的反驳，布莱希特以其独特的方式和辩证法家的态度，证明基于“普遍”规律的辩证法理论于历史实践是不利的，甚至是荒唐的。他写道：“如果人们向这种简单然而动人的观点的信徒指出，他们的看法就像那些看手相的人的看法一样，认为在看完手相确定即将到来的事情之后，能够阻止这些事件的发生，那么这些信徒就会郁郁不乐，咕哝发火。”柏拉图对话中那些无所不知、自以为是的智术师被他们自己的武器打败以后，也是这般束手无策。实际上，辩证法并不是普遍的、包含了自然和历史的规律科学，而是如布莱希特所言：“一种思维方法，或者不如说是一系列彼此相连的可理解的方法，它们允许将一些僵化的观念消解掉，用实践去反对处于统治地位的意识形态。”\footnote{布莱希特，《政治与社会文集》第1卷，第251页。}

我想要证明，黑格尔的进步理论有别于后来的进步的意识形态，它是要使辩证法的这种双重观点生效。

\section*{2}

在我们着手证明之前，似乎有必要首先简述一下黑格尔历史哲学赖以产生的历史前提。现代进步概念的历史始于16世纪，一开始是不停地发展，然后到启蒙运动时期导致了一系列二律背反。这些二律背反构成了黑格尔为此概念进行辩证奠基的起点和出发点。现代进步观念起源于城市商业和制造业等级的解放理想。随着培根、笛卡尔、霍布斯和莱布尼茨对科学理论和行动理论提出新的筹划，现代进步概念也就产生出来。这种新的理论不再像经院主义—亚里士多德哲学那样将它们的目标预设在一个最高的善或者最高的知识中，而是认为它们自身原则上是开放的，而不是封闭的。自然科学构成了这种理论的典范，其整齐划一的直线运动的概念预设了自然界并无固定的目的和目标，也预设了运动相对于静止的优先性。正如自然的科学一样，其实践上的推论，即通过劳动和行动控制自然，也不是完美的、自身封闭的，而是一个迈向未来的过程。\footnote{参见培根：《新工具》第一卷，1620年，80-84。}现在，科学和行动的过程性成为伦理学和人类学的基本概念。霍布斯用一种善(bonum)代替了旧道德哲学著作中的至善(summum bonum)，它不像后者那样在于完满和静止，而在于无休无止地从满足愿望到满足愿望的行进。\footnote{霍布斯：《利维坦》（1651年），第一部分，第11章。}近代市民社会于其产生时即被理解，它的至善包含了它对幸福的解放性要求。莱布尼茨表达了同样的观念，他理解的幸福不是尽善尽美的快乐，而是向新的完善和新的快乐不断迈进。\footnote{莱布尼茨：《自然和恩典的理性原则》(Vernunftprinzipien der Natur und der Gnade, 1714)，第18节。}这个无限进步的概念成为18世纪启蒙运动理解人的自然和人的历史的关键。莱布尼茨引入的关键词叫可完善性，它用自己的手段摧毁了传统的自然目的论，并在本体论上使市民阶级的解放要求合法化。按照这种将进步提高到宇宙论层次的理论，事物并不因其符合亚里士多德的实体和本质的本体论要求永恒不变就是完美的，而是相反：事物的可变性构成它们在一种更高意义上可以称之为完美的根据。完善只有作为可完善性（也就是说，作为完善的过程），才是可能化。

这里，进步思想达到了其最宽广的范围，由此产生了那些贯穿了启蒙运动历史哲学思想的二律背反。第一个二律背反是在【进步】概念自身中的二律背反，这个概念一方面似乎包含了一个目的——万物的完善，等等，然而另一方面却将此目的表现为原则上是不可达到的。由此产生了第二个二律背反，它是在运用这个概念时产生的。换言之，我们要么假定进步只对个人有意义，要么想要反过来证明，只有类、人类才能达到完善的目标。这个二律背反在18世纪文学上、在门德尔松和莱辛争论中得到了解决。\footnote{参见门德尔松《耶路撒冷》，《全集》第5卷，莱比锡，1844年，第120页；莱辛，《人类的教育》，《全集》第8卷，柏林，1956年，第590页以下。}第三个二律背反形式上看似乎重复了第二个二律背反的立场，但内容上却存在着原则性差异。它接受了曼德维尔和卢梭对市民社会经济、文化和政治矛盾的发现，由此，在18世纪下半叶，解放的理想便进到了批判阶段：社会可完善性的原因不是德性，而是恶德，其结果不是平等，而是不平等，其前提不是贫富均衡，而是贫富对立。早期启蒙的进步乐观主义证明其自身是一个为了保证概念自身得到应用而必须抛弃的幻想。这出现在卢梭这里。在他看来，所有超出人类自然状态的进步，既是人类文化和理性完善的脚步，也同样是人类的衰落。\footnote{《论人类不平等的起源》，《卢梭选集》，巴黎，无出版年份，第65页。}康德得出了正好相反的结论：对人类来说，人走出自然状态的过程，从单纯动物的野蛮状态过渡到人道状态，乃是从坏到好的进步，对个人来说，则并非如此。\footnote{康德，《人类历史起源臆测》（1786年），对第一部分的解说。（“从坏到好的进步”里德尔原文为“Fortschritt vom Bessern zum Schlechteren”，兹据康德原文更正。康德原文见Kants Werke. Band VIII. Akademie-Textausgabe, Berlin 1968, p. 115。——译者）}

第三个二律背反的决定性观点是由卢梭和康德阐述的一种进步辩证法的观念，这种观念使他们二人发现，历史的运动结构自身就是辩证的，是一种进步和回归的运动。这种观点构成了黑格尔的直接前提。然而卢梭对于在市民社会及其文化的基础上解决进步的辩证法仍然心存疑虑，并且着眼于一种依然与自然相契合的社会形式，筹划其《社会契约论》的模式，康德则将进步辩证法作为其历史哲学的出发点。这种历史哲学不再是将此领域以前的努力变成空洞无益的游戏的种种虚构。在特别严肃的卢梭那里，历史根本而言是一部哲理小说，因为回归的运动与退步的运动并无本质区别。康德用一种进步理论回答卢梭，认为进步就是人类所有原始禀赋的发展，并将通过经验后天认识到的现代市民社会的对抗性——人的非社会的社会性——提升为历史运动的动力，同时又按照一条先天的线索，将其提升为“人类合法秩序的原因”。通过证明人类行为中的有规律的进程及其可能同意组建一个社会（在这个社会中，每个人的自由与所有人的自由都能按照一条普遍法则而共存），解放的模式便得到了辩护。康德在其《世界公民观点之下的普遍历史观念》（1784年）第五个命题中说：“自然迫使人要去解决的人类的最大问题，就是建立一个普遍的照管权利的市民社会。”\footnote{康德原文见Kants Werke, Band VIII, p. 22。（中译文参见〔德〕康德：《历史理性批判文集》，何兆武译，商务印书馆1990年版，第8页。——译者）}康德历史建构中的先天线索，就是在权利思想中得到保证的自由，而不是至善（或者德性和幸福）的实现。历史哲学中的这个哥白尼式转向在于，偏离旧的神学和伦理-教育的历史思辨，转向现代市民社会及其历史，现在，这个历史需要一个新的、非形而上学的神正论。康德高呼：“因为在没有理性的自然界之中我们赞美造化的光荣与智慧并且把它引入我们的思考，究竟又有什么用处呢；假如在至高无上的智慧的伟大舞台上，包含有其全部的目的在内的那一部分——亦即全人类的历史——竟然始终不外是一场不断地和它相反的抗议的话？这样一种看法就会迫使我们不得不满怀委屈地把我们的视线从它的身上转移开来，并且当我们在其中永远也找不到一个完全合理的目标而告绝望的时候，就会引导我们去希望它只能是在另外一个世界里了。”\footnote{《康德全集》（科学院版）第8卷，第30页。（中译文引自《历史理性批判文集》，第20页。——译者）}

\section*{3}

黑格尔的历史哲学由此出发，我的主张是，只有它才完成了对哥白尼式转向后果的思考，并将18世纪进步理论的那些二律背反抛到了身后。黑格尔同意康德，认为世界历史的基础不是幸福或者德性与人道的教育，而是在自由与权利关系之中的自由。黑格尔称莱辛和赫尔德的“人类教育”观念是“富于精神的”，但是这种观念“只是远远地”触动他，并非“此处所谈”。\footnote{黑格尔：《历史中的理性》，第150页。参见《法哲学原理》，第343节，附释，《全集》（百年纪念版）第7卷，第447页以下。}他与卢梭、康德一致，认为可完善性概念本身无法满足，在应用到历史中时是没有标准的：“在本讲座中，一般说来，进步具有量的形式。更多的知识，更好的教养，只能是这样的比较；可以一直这样说下去，而不言及任何规定性，说到某种质的东西。”\footnote{黑格尔：《历史中的理性》，第150页。}然而，比一致更加重要的是在康德和黑格尔的历史进步观之间存在的区别。我将其总结为五点，其中第一点是决定性差异，它决定了其余几点。

1. 不同于康德及其18世纪前人，在黑格尔看来，一个普遍的照管权利的市民社会的理念不再是悬设，而是现代世界的现实性，也就是那个在法国革命中作为世界历史的原则而使所有人的自由都变得有效并得到承认的“宪政”国家的原则。在它面前，作为历史目标（同时历史也从来就不能达到这一目的）的进步的二律背反已然失效。

2. 在康德这里，进步的过程有两个要件，一是在历史之前即已存在于其最终形式之中的自然的计划，二是带有作为质料的激情的人的自然，在这种质料中，前述形式得以实现其自身。\footnote{我在这里使用了科林伍德的一个中肯的表达，见R. G. 科林伍德：《历史哲学》，斯图加特，1955年，第119页。}黑格尔则将进步置于历史自身当中，由此，历史即作为自由意识的进步而得到理解。在黑格尔看来，一般而言，历史哲学的问题并不在于能否看出自然或者神圣天意的计划，或是将这种计划作为外在地考察经验历史的先天线索，而是在于一次性的历史过程是否包含了一种为其辩护的意义——这个问题的解答存在于时间和概念的辩证法当中，对此下面还会言及。

3. 与康德及其18世纪前人相对，在黑格尔这里，在理性支配的未来与无理性的过去之间并不存在一道鸿沟，因为对他来说，进步的目标不在未来，而在当下。\footnote{参见上书，第111页。}从当下出发，可以将历史回头构建为一个连续向前发展的过程。可以说，这里有着黑格尔历史哲学的限度，但也是其伟大之所在：通过历史连续性的思想，将历史中的进步合法化了。

4. 在康德及其前人那里，作为进步的历史活动于个人与类的关系的二律背反之中，但在黑格尔看来，这种关系就其本身及其与历史的关联而言都只是一种抽象物。就其本身而言是抽象的，因为在个体的自知、它的自我意识(Ich-Bewusstsein)中，类是同时得到思考的，并且其中也就有“我即我们”，或是黑格尔以“精神”之名为解释历史的进步奠定基础的那种原始的综合性的总体意识，在此总体意识中，个体和类两者达到了自由法意识；在与历史的关联中是抽象的，因为个体在其要素中本质上就是“普遍的自然”，即一个特定的民族，而类是一种严格意义上的历史的类，也就是那种世界精神的形态，它从各个民族和国家的过程中作为它们的历史而产生出来。\footnote{参见黑格尔，《法哲学原理》，第340节，第446页。《哲学全书》，第548节，《全集》（百年纪念版）第10卷，第426页以下。}

5. 尽管康德在回答“人类是否应该理解为不断地朝向改善而进步”时，\footnote{参见康德，《系科之争》（1798年），第二篇，1—10。}为了未来而先天地考虑了历史的问题——据此，讲史者制造并安排了他事先预告的那些历史事件——但他并未为了当下和过去而考虑存在于这种观念中的实践观点、人的创造性活动在历史中的作用：在一定程度上，实践被自然的计划消解掉了。相反，尽管黑格尔这里也有与康德的自然计划类似的“理性的狡计”学说，但他将历史、精神的王国视为人类劳动的产物和结果：“精神的领域由人产生。人能对神的国度提出各种各样的见解，因此总有一个精神的王国在人世中得到实现，并被人设定于实存之中。”\footnote{黑格尔：《历史中的理性》，第50页。}在字面意义上，历史中的进步由人所拟就：“然而，事实上，我们之所是，同时就是历史的【……】属于我们的、属于当今世界的理性的财富，不是直接产生的，也不是成长于当下的土壤；相反，其本质在于，它是一种遗产，进一步说，它是劳动的结果，并且是人类历代祖祖辈辈劳动的结果。”\footnote{黑格尔：《哲学史讲演录》第1卷，《全集》(百年纪念版)第17卷，第28页。}

\section*{4}

这里，我不详细讨论第二至五点，而只论述第一点。它最好地说明了我所言的“黑格尔是康德引入历史哲学中的哥白尼式转向的完成者”的主张。为此目的，我对康德和黑格尔与法国大革命之间的关系进行了比较。这个比较不是外在的表态，而是要突出这场革命对于他们的历史哲学思想的价值和意义。康德在18世纪90年代开始将人类统一进步的理念看作自然的单纯意图，以及在一定程度上与历史过程并列的想象的历史目的，并且已经致力于将人类统一进步的理念与历史过程联系起来。与《纯粹理性批判》中的图式论问题——应当思考单纯的知性概念如何运用于可能的经验对象——相类似，现在历史哲学也产生了在理念中设想的、通过自然意图予以保证的历史目的与实际的历史过程之间的关系问题。

康德用历史符号理论解决这个理念在经验中的图式化问题。这个问题在于，找出一个历史事件，它在时间上是不确定的，却暗示了人类作为其向改善推进原因的特性和能力，由此即可得出到达那种历史目标之可能性的结论——“然后这一结论还要能够这样地扩大到以往时代的历史（即它永远是在前进的），以至于那个事件的本身并不必须被看成是这种历史的原因，而是必须被看成只不过是一种示意、是一种历史符号(Signum rememorativum, demonstrativum, prognostikon)【……】”。\footnote{康德，《系科之争》，第二篇，5。（中译文引自〔德〕康德，《论教育学》，赵鹏、何兆武译，上海人民出版社2005年版，第116页。——译者）}对康德来说，这种历史符号——这里他的进步理论的二律背反又重新开始了——不是法国大革命，而是虽未参加然而同情法国革命的公众的道德反应。黑格尔在历史哲学讲座中报道了法国革命那些观众的热情和崇高情感；但在他看来，真正的历史符号乃是法国革命事件本身，其中包含了一种对于以后和以往的意义：对于以后具有意义，是因为现在，权利和自由的思想在历史上第一次被提高为政治制度的基础，“从此以后，一切都应该建立在这个基础之上”；对于以往具有意义，是因为自从阿那克萨哥拉将理性提高为宇宙的原则以来，哲学（再次）出现在这种思想的符号中。同时，这是从理性作为宇宙原则到理性作为历史原则的何种飞跃和迈进啊！黑格尔这样说明这种进展：“自从日丽中天、众星绕它运行以来，还从未见过此种事情：人立于头脑（即思想）上面，按照思想构筑现实。阿那克萨哥拉最早说过，努斯($\nu o\tilde{\upsilon}\varsigma$)统治世界；然而直到现在，人们才认识到，思想应该统治精神的现实。”\footnote{黑格尔：《世界历史哲学讲演录》第4卷，G. 拉松出版，莱比锡，1920年，第926页。}黑格尔以这种方式将法国革命本身理解为历史符号和权利与自由思想的现实化，由此，康德进步理念的图式化问题便获得了不同的解答。黑格尔用时间与概念、自我意识与时间的辩证法取代了那种图式化，由此，他的思想便实现了向现代世界以及发生在现代世界基础上的历史过程的突破。

哲学中揭橥自由之基础的事件，作为时间中的事件，发生在法国革命之前。卢梭在自由意志的思想中发现了一切权利的实体基础，康德在“我思”、意识的先验统一中发现了一切经验的最高条件和所有思想规定的起源。这两个事件共同构成了黑格尔哲学的基本经验。黑格尔哲学从与自身同一的自我的思想出发，同时结合费希特和早期谢林的思想对其进行了解释。自我不是空洞的统一性；相反，按照概念的系统构造，它是三个自身区分开来的环节的统一。首先，在康德的“我思”的意义上，是纯粹自身与自身发生关联的统一性，通过抽除所有内容，这种统一性回到与自身的无限制等同的自由之中，并将自己规定为普遍性；其次，是绝对规定了的存在，此种存在与他者对立，并排斥他者，此即个别性或者个人的人格性；第三，是绝对的普遍性，它同时又是直接的绝对个别化——自我的具体结构，黑格尔将它与概念的结构等同起来：普遍物与个体合而为一。\footnote{黑格尔：《逻辑学》，第二部分，莱比锡，1951年，第219页。}

“精神”就是这个结构向——黑格尔用来说明社会历史世界及其变化的——那个总体意识的辩证展开。然而，并不是“精神”的概念才包含了与时间的联系；相反，概念的概念、自我亦然——这是超出康德以及之前所有哲学的决定性步伐。在康德那里，尽管“我思”必然伴随时间性的表象过程，但它自身却在时间之外；对黑格尔来说，自我根本上就在时间之中：“自我在时间之中，时间就是主体自身的存在。”\footnote{黑格尔，《美学》，第三部分，《全集》（百年纪念版）第14卷，第151页。}自我不是从它与时间的关系出发规定自身；相反，它从时间与自我或概念的关系出发规定时间：时间就是“定在着的概念自身”\footnote{黑格尔，《哲学全书》，第258节，注释，《全集》（百年纪念版）第9卷，第258页；《精神现象学》，《全集》（百年纪念版）第2卷，第44、612页。}，或如黑格尔所言，时间“与自我（即纯粹自我意识的自我）是同一个原则”。在自我与时间的辩证法中，黑格尔实现了体系向历史的突破，一笔勾销了康德的图式化问题。康德则还想要用图式化去克服独断论形而上学的隐秘神话学以及与之相应的道德恐怖主义（人类发现自己不断朝着更恶倒退）、幸福主义（人类不断朝向改善前进）和阿布德拉主义（人类发现自己永远处于静止状态，愚蠢地忙忙碌碌乃是人类的特性）的荒谬历史主张。\footnote{参见康德，《系科之争》，第二篇，3，a-c。}

自我与时间的辩证关系使得在历史中指出康德将其适用于自然经验的那种联系和统一性成为可能。自我并不是时间在上面随意涂画的白板，也不再如那些独断的历史观念所见（康德的历史构想部分地也是如此），是那些相互交错的事件和意见的漠不相关的、没有自我的中途点\footnote{参见狄尔泰，《施莱尔马赫传》(Das Leben Schleiermachers)第2卷，2，哥廷根/柏林，1966年，第538页。}，而是时间和历史的统一，并在这种统一中自己出场。黑格尔把自我的这种出场叫作自我的自由。时间与自我（即纯粹自我意识的自我）是同一个原则，而纯粹自我意识同时就是自由的原则。\footnote{黑格尔：《哲学全书》，第258节，注释，《全集》（百年纪念版）第9卷，第79页。}

黑格尔按照自由意识进步理念的线索，基于这种自我、时间和自由的历史辩证法，对历史进行解释。这种辩证法将黑格尔哲学的基本经验融入其哲学体系之中。现在，在法国革命的历史事件中得到实现的市民社会的理性理想，就能通过实际发生的历史而得到辩护。我在其中看到了由康德引入的历史哲学中的哥白尼式转向的完成。这种通过时间性事件所做的辩护并不意味着，当黑格尔在前后相继的世界事件的尺度上按照成果理性对进步进行量度时，将历史化为一个线性的过程。情况恰好相反。当黑格尔预设的市民社会的理性理想在国家中成为现实，而现代世界的社会历史现实变得理性时，则自由的意识也就不再归结为单义的进步运动，而是归结为双义的进步和回归的运动。悖论的表达就是：真正的进步、它的完成就是回归——黑格尔的理论用它实现了精神的体系概念，即精神在它自己的别处也是自在的或是自由的：“理性认识真实的东西，自在自为的存在者，这种东西是没有限制的。精神的概念就是回归自身，以自己为对象；因此，前进就不是所向无定地向无限物进发，而是有一个目的，即回归自身。”\footnote{黑格尔：《历史中的理性》，第181页。参见同书，第58、65、67、70页以下。}回归的结构既不是说，进步在历史中到达了终点——黑格尔排除了这种看法，因为从地理上看，自由的实现，作为世界历史的目标，仅仅出现世界的一个角落，并且“从昨天”才开始\footnote{参见黑格尔：《法哲学原理》，第62节，附释。}，同时它也没有预设，历史运动就是重复——同时代人的复辟企图意义上的或是更加古老的、直到维科为止一直有效的政治制度的循环论意义上的——过去的社会形式。黑格尔拒斥浪漫派国家理论中的中世纪教义，称之为“我们时代的无聊，想把其中的精髓变成口号”。\footnote{黑格尔：《世界历史哲学》第4卷，第840页。}尽管他想要在他自己的时代和晚期罗马帝制时代之间确定一些共同的结构，但他绝不认为在历史衰落时期之间具有“重大相似性”，而在从福尔格拉夫(Karl Vollgraff)到拉绍尔克斯(Ernst von Lasaulx)和布克哈特再到斯宾格勒的19世纪保守派历史理论中，这种相似性发挥了显著作用。

这样，内在于历史运动之中的进步，就像历史运动一样，是不可逆的。这并不是说，进步就是直线式的，以一种直接的、无对立的自然过程的方式行进的。如同历史中的所有变化一样，每一进步都是矛盾的环节——有和无、产生和消逝、进展和衰落——的统一。至关重要的是，黑格尔将变化思考为一个缓慢的、没有断裂的过程：“用变化的渐进性来理解产生和消逝，就是同语反复所特有的无聊；那意味着：正在产生或者消逝的东西，预先就已经是现成的了，而变化则成了外在区别的简单改变，这样实际上就是同语反复。”\footnote{黑格尔：《逻辑学》，《全集》（百年纪念版）第4卷，第461页。（中译文引自《逻辑学》上卷，第405页。译文略有改动。——译者）参见伽兰第(R. Garaudy)：《上帝死了。对黑格尔的一个研究》(Gott ist tot. Eine Studien über Hegel)，柏林，1965年，第386页以下。}因此，进步首先是在历史运动的关节点和转折点上、在时间与自我发生分离的时期得到实现的。在这种时期，一个个人或一个民族的精神不再与其制度相一致，制度也不再与精神一致。于此，可以区分出进步的三种模式，据黑格尔看来，它们都在历史中出现，同样必然，同样正当：革命的模式、改革的模式以及最后作为世界精神的终极理由(ultima ratio)——一个民族被排除在历史进步民族的序列之外：“如果一个国家的法制所表示的真理已经不符合于它的自在本性，那么它的意识或概念与它的现实性就存在着差别，它的民族精神也就是一个分裂了的存在。有两种情况可以发生：这民族或者由于一个内部的强力的爆破，粉碎了那现行有效的法律制度，或者较平静地、较缓慢地改变那现行有效的但却已不复是真的伦理、已不能表现民族精神的法律制度。或者（第三种模式即在其中）一个民族缺乏理智和力量来作这种改变，因而停留在较低劣的法律制度上；或者另一个民族完成了它的较高级的法制，因而成为一个较卓越的民族，与之相对，而前一个民族必定会不再为一个民族，并受制于这个较卓越的民族。”\footnote{黑格尔：《哲学史讲演录》第2卷，《全集》（百年纪念版）第18卷，第276页以下。（中译文引自《哲学史讲演录》第二卷，第249-250页。——译者）}第三种模式包含了作为世界历史的现实历程出现在《法哲学原理》结尾的那些有限的民族精神的那种“辩证发展现象”。这是历史哲学在哲学体系中的位置。如果我们想要在黑格尔的体系和历史统一的视角下理解康德引入的方法论转向的意义，就必须说明黑格尔历史哲学的这种位置。

\section*{5}

在黑格尔的阐述中，国家的普遍理念在其完成的顶点下降为在历史的活动中彼此对立的“四种”特殊的“国家”。现代世界的国家与市民社会处于差异的关系之中，自己平衡了市民社会的内部差异。它不再是传统的政治状态概念，也不是古典政治学和近代自然法理论的不变的和无历史的自然模式或契约模式；相反，历史——作为实践哲学的新的上司——高踞其上。在历史中，“伦理整体本身、国家的独立性都被委诸偶然性之下”（第340节）。黑格尔在《法哲学原理》体系最后引入的历史，就是从霍布斯到卢梭的自然法学说于其体系开端处设定的那种自然状态理念的现实性。在黑格尔这里，自然法理论家们设想的从自然状态到市民社会并在与国家的同一中静止下来的运动，不是始于国家同市民社会的关系，而是始于与其他国家的关系。这种自然状态不是虚构的状态，而是现实的状态，是历史的运动。《法哲学原理》将这个运动接受到其体系之中，从而摆脱了法和社会的抽象自然理论。在世界历史中，“普遍精神”的定在的要素，作为法概念自我展开的起点和终点，就是“精神性的现实性”，在这种现实性中，人和物的“自然”就不是固定不变的规律，而是黑格尔思考为自由的概念，它“在其内在性和外在性的全部范围内”实现其自身。在这个与法哲学体系结束和历史哲学开端同时发生的运动的辩证法中，时间、概念与自由统一起来了。历史作为自由意识中的进步，就是对于在现代宪法国家中得到实现的那种自由的辩护。

由此出发，即可明白黑格尔历史哲学提出的非同寻常的要求。历史哲学的主题和对象不再是对历史过程本身无甚兴趣、根据一条先天的线索写就的一种可能的历史的理念，而是可以理解为实在性的历史本身。这种要求并非现在才被人们斥为幻想，而我也绝非想要以一个正统黑格尔派的态度着力为其辩护。最后，在我为其辩护时，也只想提出一点意见，这首先是为了质疑那些罕见一致地臆想自己正在接近一个后历史时代的现代学派——结构主义、实证主义和精神科学的解释学——对解放的批判的历史概念的抹黑。事实上，黑格尔的辩证历史解释提出的要求可以简化为列维—斯特劳斯\footnote{参见列维—斯特劳斯：《野性的思维》，巴黎，1962年，德文版，美茵河畔法兰克福，1968年，第292页。}拿来解释萨特的辩证理性批判问题的那个表达式：在何种条件下法国革命的神话是可能的？但是人们可以从黑格尔哲学得到一个标准（这在萨特这里也许没那么容易），从而不可能将历史和神话混同起来。并且我认为，这个标准存在于黑格尔确定的概念（或自我）与时间的关系中。正如时间的结构必然以自我的结构为前提，自我也同样以时间为前提：在严格意义上，革命是时间中的事件，是历史符号，它作为这样的符号意味着进步，而不是重复同一结构。古典政治（我用它表示其哲学维度。在我看来，在此维度中，必然要讨论神话和历史问题）要么回溯到神话（如柏拉图从《理想国》到《蒂迈欧》的宇宙论及其关于时间的有规律的轮回学说的过渡），要么回溯到循环的时间结构的神话考古学（如亚里士多德《政治学》第1卷），以巩固城邦的存在；现代自然法则在源自历史的自然状态的学说中进行虚构，创造出一部哲理小说，以使国家作为市民社会的生成过程能够止步于国家之中——已经有了一部历史，然而再也没有历史了。黑格尔让他的《法哲学原理》体系纲要结束于对市民社会和国家之实在性的历史展望中，在这种实在性面前，市民社会和国家都“只是作为观念物而存在”。\footnote{黑格尔，《法哲学原理》，《全集》（百年纪念版）第7卷，第446页。（原引文引自《法哲学原理》第341节，而非第340节。中译文参见《法哲学原理》，第351页。——译者）}出现于历史中的辩证法确实是一种辩护的辩证法，但也是一种克服的辩证法。黑格尔说过，唯有当下存在，以前和以后都不存在；\footnote{黑格尔，《哲学全书》，第二部分，第259节，补充，《全集》（百年纪念版）第9卷，第86页。}然而他还补充道，具体的当下是过去的结果，并且孕育了未来。真正的当下是概念把握到的时间整体。